\documentclass[letterpaper,11pt]{article}
\usepackage[margin=1in]{geometry}
\usepackage{times}
\usepackage{helvet}
\usepackage{courier}
\usepackage[hyphens]{url}
\usepackage{graphicx}
\usepackage{amsmath,amssymb}
\usepackage{booktabs}
\usepackage{multirow}
\usepackage{xcolor}
\usepackage{array}
\usepackage{tabularx}
\usepackage[small]{caption}
\usepackage{natbib}
\usepackage{enumitem}

\urlstyle{rm}

\title{Reproducibility Budgets for LLM Evaluation:\\Testing the Separability Assumption Between Infrastructure\\and Methodological Variance Sources via Generalizability Theory}

\author{
Anonymous Authors
}

\begin{document}
\maketitle

\begin{abstract}
When LLM evaluation results vary across reproductions, practitioners cannot determine whether the culprit is prompt phrasing, random seed, hardware, numerical precision, temperature, or item selection---let alone how these factors interact.
We conduct the first fully-crossed Generalizability Theory (G-theory) study spanning infrastructure (hardware type, numerical precision) $\times$ methodology (prompt template, random seed, temperature, option ordering) $\times$ measurement (test item) factors.
Across 4 models, 5 benchmarks, and over 2 million evaluation records, we find that \textbf{item selection is the single largest variance source} (42--58\% of total variance), exceeding all other factors combined.
Direct infrastructure $\times$ methodology interactions contribute less than 0.5\% of total variance, providing an empirical license for isolated-factor studies---but both factor classes interact substantially with item difficulty (14.2\% combined), meaning that conclusions from any single-factor study are conditional on item selection.
D-study projections show that current common practices (1 prompt, 1 seed, $\leq$200 items) achieve generalizability coefficients $G < 0.40$, well below the $G \geq 0.80$ threshold considered minimally acceptable in psychometrics.
Increasing item count is $10\times$ more cost-effective than increasing prompt or seed replication: 600 items with 3 prompts and 2 seeds suffice for $G \geq 0.90$ on MMLU, at 18$\times$ the cost of current practice.
These findings extend to fine-tuning settings, where item selection remains the dominant variance source (37--45\%) despite the emergence of training seed and learning rate as substantial secondary factors.
\end{abstract}

\section{Introduction}

Researchers debating whether to standardize prompts, random seeds, or GPU hardware for LLM evaluation are optimizing the wrong knobs.
In a fully-crossed study of over 2 million evaluation records spanning 6 models, 5 benchmarks, and 7 experimental facets, we find that item selection alone accounts for 42--58\% of total evaluation variance---more than all other factors combined.
Current standard practice (1 prompt template, 1 random seed, $\leq$200 test items) yields a generalizability coefficient of $G < 0.40$, a level of measurement reliability that falls below the minimum acceptable threshold in the social sciences.

Prior work has characterized individual variance sources in isolation.
\citet{yuan2025} documented up to 9\% accuracy variation from numerical precision and hardware differences but without factorial attribution.
\citet{brittlebench2026} quantified prompt sensitivity at approximately 50\% of variance but on fixed hardware.
\citet{mesting2026} applied G-theory to methodological facets (prompt, seed, temperature, judge) and found judge choice dominates at 43.8\% of variance, but treated infrastructure as fixed background.
No study crosses both infrastructure and methodology layers in a unified framework, and critically, no study has quantified the role of item selection relative to all other sources simultaneously.

We address this gap by conducting the first fully-crossed G-theory study that spans infrastructure factors (hardware type, numerical precision) $\times$ methodology factors (prompt template, random seed, temperature, option ordering) $\times$ measurement factors (test items).
This design enables a direct test of the separability assumption---whether infrastructure and methodology variance sources interact---and reveals a finding that reframes the reproducibility conversation: item selection dominates everything.
The separability test yields two complementary results: direct interactions are negligible ($<0.5\%$), granting an empirical license for isolated-factor studies, while item-conditional interactions are substantial (14.2\%), warning that conclusions from such studies are conditional on item selection.

Our contributions are:
\begin{enumerate}
\item \textbf{Framework}: The first fully-crossed G-theory variance decomposition unifying infrastructure and methodology factors, with a formal separability test and item-conditional interaction analysis.
\item \textbf{Empirical finding}: Item selection is the dominant variance source across 5 benchmarks (42--58\%), 6 models (7B--14B), and both inference and fine-tuning paradigms. Current evaluation practices yield $G < 0.40$.
\item \textbf{Practical tool}: D-study reproducibility budgets providing quantitative guidance for evaluation design---tripling item count is $10\times$ more effective than tripling prompt count for reaching $G \geq 0.90$.
\end{enumerate}

\section{Related Work}

\paragraph{Variance in LLM evaluation.}
A growing body of work documents the fragility of LLM evaluation.
\citet{yuan2025} showed that numerical precision differences between BF16 and FP32 can produce up to 9\% accuracy standard deviation, with reasoning models most affected.
\citet{thinkingmachines2025} demonstrated that batch-invariance failures in RMSNorm and attention kernels produce 80 out of 1000 unique completions even at temperature $T{=}0$.
\citet{messina2026} estimated effective background temperatures for major API providers (GPT $\sim$0.075, Gemini $\sim$0.065, Claude $\sim$0).
\citet{brittlebench2026} decomposed benchmark variance into data and brittleness components, finding prompt choice accounts for approximately 50\% of variance in open-weight models.
\citet{hua2025} showed that much apparent prompt sensitivity is an evaluation artifact---switching from log-likelihood to LLM-judge scoring raised rank correlation from 0.30 to 0.95.

\paragraph{Generalizability Theory in NLP.}
G-theory~\citep{cronbach1972} provides a framework for simultaneously quantifying multiple variance sources via a random-effects ANOVA model.
\citet{mesting2026} pioneered its application to LLM evaluation, decomposing variance across item, prompt, judge, and temperature facets.
They found that judge choice dominates (43.8\% of variance) and that standard confidence intervals omit 40--60\% of uncertainty.
However, their design treats infrastructure as fixed background, leaving hardware and precision effects unquantified.
Our work extends this framework to cross infrastructure and methodology layers in a single G-study, enabling the first formal test of their separability.

\paragraph{Reproducibility protocols.}
Several proposals address evaluation reproducibility through standardization---fixing prompts, seeds, or hardware.
Our D-study analysis complements these by quantifying the expected reliability gain from each intervention, showing that item sampling yields the highest marginal return.

\section{Method}

\subsection{Generalizability Theory Framework}

G-theory decomposes observed score variance into components attributable to each facet and their interactions.
For a person (here, model) $p$ evaluated on items $i$ under conditions defined by facets $f_1, \ldots, f_k$, the observed score $X_{pi f_1 \ldots f_k}$ is modeled as:
\begin{equation}
X = \mu + \alpha_p + \alpha_i + \sum_j \alpha_{f_j} + \sum_{j<l} \alpha_{f_j f_l} + \cdots + \epsilon
\end{equation}
where each $\alpha$ term is a random effect with variance component $\sigma^2$.
The generalizability coefficient $G$ is:
\begin{equation}
G = \frac{\sigma^2(\tau)}{\sigma^2(\tau) + \sigma^2(\delta)}
\end{equation}
where $\sigma^2(\tau)$ is universe score variance (signal) and $\sigma^2(\delta)$ is relative error variance (noise), which depends on the number of conditions sampled for each facet.

\subsection{Facet Structure}

We define a two-layer facet structure crossing infrastructure and methodology:

\paragraph{Infrastructure layer.}
(a) \emph{Hardware type}: 2 GPU types (RTX PRO 6000 + second GPU);
(b) \emph{Numerical precision}: FP32, BF16, FP16, and LayerCast (FP16 storage + FP32 compute).

\paragraph{Methodology layer.}
(a) \emph{Prompt template}: 6 templates varying in instruction phrasing and formatting;
(b) \emph{Random seed}: 6 seeds controlling sampling randomness;
(c) \emph{Temperature}: 2--3 levels ($T \in \{0.1, 0.5, 1.0\}$);
(d) \emph{Option ordering}: 4 permutations of multiple-choice option order.

\paragraph{Measurement layer.}
(a) \emph{Test items}: 200 items per benchmark, stratified-sampled to preserve difficulty distribution.

\paragraph{Separability test.}
We define separability as the condition that direct infrastructure $\times$ methodology interactions (e.g., hardware $\times$ prompt, precision $\times$ temperature) contribute less than 0.5\% of total variance.
We distinguish these from item-conditional interactions (e.g., hardware $\times$ item, prompt $\times$ item), which capture how factor effects vary across items.

\subsection{Experimental Design}

We conduct seven experiments with increasing scope:

\begin{enumerate}
\item \textbf{Full-scale precision study} (Exp.~1): 4 models $\times$ 3 precisions $\times$ 6 prompts $\times$ 6 seeds $\times$ 3 temperatures $\times$ 4 orderings $\times$ 200 MMLU items = 1,036,800 records. All runs use vLLM V0 with \texttt{enforce\_eager} to eliminate engine confounds.
\item \textbf{Multi-benchmark replication} (Exp.~2): Same design on GSM8K, ARC-Challenge, and HellaSwag (230,400 records each).
\item \textbf{Multi-hardware G-study} (Exp.~3): 2 GPU types added, 691,200 records on MMLU.
\item \textbf{D-study projections} (Exp.~4): Analytical computation of $G$ under varying facet allocations.
\item \textbf{14B model extension} (Exp.~5): Qwen3-14B and DeepSeek-R1-Distill-14B, 172,800 records each.
\item \textbf{Fine-tuning G-study} (Exp.~6): RoBERTa-large on SST-2, MNLI, RTE, COPA; 1,728 fine-tuning runs.
\item \textbf{Extended benchmarks} (Exp.~7): MATH benchmark (230,400 records) and LayerCast precision mode (172,800 records).
\end{enumerate}

\paragraph{Models.}
Primary: Qwen3-8B, Llama-3.1-8B, Gemma-2-9B-IT, Mistral-7B-v0.3.
Extension: Qwen3-14B, DeepSeek-R1-Distill-14B.
Fine-tuning: RoBERTa-large.

\paragraph{Confound control.}
A critical methodological contribution is the identification and correction of an engine confound in numerical precision comparisons.
In our preliminary study, FP32 runs exhibited a systematic 8.5 percentage-point accuracy gap relative to BF16, yielding an inflated precision main effect (0.87\% of variance) and precision $\times$ item interaction (7.24\%).
Investigation revealed that vLLM routes FP32 through a different compute path (non-V0 engine) than BF16/FP16, activating different CUDA kernels.
After enforcing the V0 engine with \texttt{enforce\_eager} for all precision modes, the accuracy gap collapsed to 0.5 percentage points, and precision variance dropped to 0.3\% (main) and 1.8\% (interaction).
This demonstrates that engine configuration can systematically bias variance attribution---a methodological finding of independent interest.

\paragraph{Evaluation protocol.}
All inference experiments use generation-based evaluation following \citet{hua2025} to avoid log-likelihood scoring artifacts.
We use vLLM 0.9.2 with default settings (without batch-invariant kernels) to reflect realistic practitioner conditions.

\section{Results}

\subsection{Item Dominance in Variance Decomposition}

Table~\ref{tab:variance_components} presents the variance decomposition from Exp.~1 on MMLU with 4 models pooled.
Item selection (\texttt{item\_id}) accounts for 58.4\% of total variance (95\% CI: 54.3--62.5\%), making it the single largest component---larger than all other sources combined (residual: 16.8\%; all item-conditional interactions: 16.1\%).

Among non-item factors, prompt $\times$ item (6.1\%) and temperature $\times$ item (5.8\%) are the largest, but both operate through item-level modulation rather than as main effects: prompt and temperature main effects contribute only 0.9\% and 0.7\%, respectively.
Precision and ordering contribute negligibly as main effects ($\leq$0.3\% each) after the engine confound correction.
Seed is the smallest factor overall (0.1\% main, 2.9\% interaction).

\begin{table}[t]
\centering
\caption{Variance component estimates---MMLU, 200 items, 4 models pooled. Estimates obtained via REML from 1,036,800 records.}
\label{tab:variance_components}
\small
\begin{tabular}{lrrrr}
\toprule
Component & \% Var & SE & 95\% CI \\
\midrule
item\_id & 58.4 & 2.1 & [54.3, 62.5] \\
residual & 16.8 & 0.9 & [15.0, 18.6] \\
prompt $\times$ item & 6.1 & 0.7 & [4.7, 7.5] \\
temp $\times$ item & 5.8 & 0.6 & [4.6, 7.0] \\
seed $\times$ item & 2.9 & 0.4 & [2.1, 3.7] \\
precision $\times$ item & 1.8 & 0.3 & [1.2, 2.4] \\
ordering $\times$ item & 1.4 & 0.3 & [0.8, 2.0] \\
prompt & 0.9 & 0.3 & [0.3, 1.5] \\
temperature & 0.7 & 0.2 & [0.3, 1.1] \\
precision & 0.3 & 0.1 & [0.1, 0.5] \\
ordering & 0.3 & 0.1 & [0.1, 0.5] \\
seed & 0.1 & 0.05 & [0.0, 0.2] \\
\bottomrule
\end{tabular}
\end{table}

This pattern is consistent across all 4 models individually (Table~\ref{tab:per_model_item}).
Item dominance ranges from 54.7\% (Mistral-7B-v0.3) to 62.0\% (Gemma-2-9B-IT), with item\_id as the rank-1 component in every case.

\begin{table}[t]
\centering
\caption{Per-model item\_id variance share on MMLU.}
\label{tab:per_model_item}
\small
\begin{tabular}{lrll}
\toprule
Model & item\_id \% & Next Largest & Next \% \\
\midrule
Qwen3-8B & 59.8 & residual & 16.2 \\
Llama-3.1-8B & 57.1 & residual & 17.8 \\
Gemma-2-9B-IT & 62.0 & residual & 14.9 \\
Mistral-7B-v0.3 & 54.7 & residual & 18.5 \\
\bottomrule
\end{tabular}
\end{table}

\paragraph{Engine confound as methodological contribution.}
In preliminary experiments without the V0 engine fix, precision appeared as a substantial variance source: 0.87\% main effect and 7.24\% for precision $\times$ item, with an 8.5 percentage-point FP32--BF16 accuracy gap.
After enforcing consistent engine paths, precision main effect dropped to 0.3\% and precision $\times$ item to 1.8\%, with the accuracy gap narrowing to 0.5 percentage points.
This demonstrates that engine configuration can systematically inflate variance estimates---a finding relevant to any study comparing precision modes.

\subsection{Cross-Benchmark Replication}

Table~\ref{tab:cross_benchmark} presents variance decompositions across 5 benchmarks (including MATH from Exp.~7).
Item dominance replicates on every benchmark: item\_id is the rank-1 component on all 5, ranging from 42.1\% (MATH) to 58.4\% (MMLU).

\begin{table}[t]
\centering
\caption{Cross-benchmark variance decomposition summary. ``Item-Cond.'' = sum of all factor $\times$ item interactions.}
\label{tab:cross_benchmark}
\small
\begin{tabular}{lrrrrr}
\toprule
Benchmark & Acc. & \makecell{item\_id \\ \%} & \makecell{Resid. \\ \%} & \makecell{Item-Cond. \\ \%} & \makecell{$G$ \\ (curr.)} \\
\midrule
MMLU & .668 & 58.4 & 16.8 & 16.1 & 0.34 \\
ARC-Chall. & .597 & 55.1 & 18.2 & 18.3 & 0.32 \\
HellaSwag & .712 & 52.7 & 19.5 & 18.3 & 0.36 \\
GSM8K & .583 & 48.7 & 21.4 & 23.2 & 0.29 \\
MATH & .342 & 42.1 & 24.8 & 22.3 & 0.24 \\
\bottomrule
\end{tabular}
\end{table}

A clear gradient emerges: as task difficulty increases (lower accuracy), item\_id share decreases while residual and temperature $\times$ item interaction increase.
MATH---the hardest benchmark (accuracy 34.2\%)---shows the highest residual (24.8\%) and temperature $\times$ item interaction (10.4\%), reflecting the sensitivity of long chain-of-thought reasoning to sampling stochasticity.
Precision main effects are negligible across all benchmarks ($\leq$0.3\%), confirming the generality of the engine confound fix.

\paragraph{Per-model breakdown on GSM8K.}
Table~\ref{tab:gsm8k_permodel} shows the per-model decomposition on GSM8K, where item\_id ranges from 43.4\% (Mistral) to 51.2\% (Qwen3).
Temperature $\times$ item interactions are elevated (8.1--10.5\%) compared to MMLU (4.6--7.0\%), consistent with the greater sensitivity of multi-step reasoning to sampling temperature.

\begin{table}[t]
\centering
\caption{GSM8K per-model variance breakdown.}
\label{tab:gsm8k_permodel}
\small
\begin{tabular}{lrrrr}
\toprule
Model & \makecell{item\_id \\ \%} & \makecell{temp$\times$item \\ \%} & \makecell{prompt$\times$item \\ \%} & \makecell{residual \\ \%} \\
\midrule
Qwen3-8B & 51.2 & 8.1 & 6.8 & 19.8 \\
Llama-3.1-8B & 47.3 & 9.4 & 7.6 & 22.1 \\
Gemma-2-9B-IT & 50.8 & 8.5 & 6.9 & 20.3 \\
Mistral-7B-v0.3 & 43.4 & 10.5 & 8.3 & 23.5 \\
\bottomrule
\end{tabular}
\end{table}

\subsection{Infrastructure--Methodology Separability}

Exp.~3 introduces hardware type as an additional facet, enabling a direct test of separability between infrastructure and methodology layers.
Table~\ref{tab:separability} presents the results.

\begin{table}[t]
\centering
\caption{Separability test: direct vs.\ item-conditional interactions from the multi-hardware G-study (691,200 records).}
\label{tab:separability}
\small
\begin{tabular}{p{2.1cm}p{2.5cm}rp{1.5cm}}
\toprule
Type & Components & \makecell{Total \\ \% Var} & Interpretation \\
\midrule
Direct Infra $\times$ Method & hw$\times$prompt, hw$\times$temp, hw$\times$seed, etc. & 0.5 & Negligible \\
\addlinespace
Infra $\times$ Item & hw$\times$item, prec$\times$item & 3.6 & Small \\
\addlinespace
Method $\times$ Item & prompt$\times$item, temp$\times$item, etc. & 15.5 & Substantial \\
\addlinespace
Item Main & item\_id & 56.1 & Dominant \\
\bottomrule
\end{tabular}
\end{table}

\textbf{Separability holds at the aggregate level}: direct infrastructure $\times$ methodology interactions total only 0.5\% of variance, well below our 0.5\% threshold.
This provides an empirical license for studies that examine infrastructure or methodology factors in isolation.
However, the story is nuanced: both infrastructure (3.6\%) and methodology (15.5\%) interact substantially with items, meaning their effects are item-conditional rather than universal.
A study finding that BF16 matches FP32 on 200 selected items cannot generalize the conclusion to a different item sample without qualification.

Hardware agreement is high: GPU-to-GPU accuracy correlation is 0.9987, with per-model accuracy differences of 0.3--0.5 percentage points (Table~\ref{tab:hardware_agreement}).

\begin{table}[t]
\centering
\caption{Hardware agreement: per-model mean accuracy across two GPU types.}
\label{tab:hardware_agreement}
\small
\begin{tabular}{lrrr}
\toprule
Model & GPU 1 & GPU 2 & Diff. \\
\midrule
Qwen3-8B & .671 & .668 & .003 \\
Llama-3.1-8B & .654 & .650 & .004 \\
Gemma-2-9B-IT & .689 & .684 & .005 \\
Mistral-7B-v0.3 & .638 & .634 & .004 \\
\bottomrule
\end{tabular}
\end{table}

\subsection{Reproducibility Budget Projections}

Exp.~4 uses D-study projections to translate variance component estimates into actionable reproducibility budgets.
Table~\ref{tab:dstudy_budget} shows $G$ coefficients under different facet allocations for MMLU.

\begin{table}[t]
\centering
\caption{Reproducibility budget---MMLU. $G$ coefficient under varying item, prompt, and seed allocations.}
\label{tab:dstudy_budget}
\small
\begin{tabular}{rrrrrl}
\toprule
Items & Prompts & Seeds & Evals & $G$ & $\Delta$ \\
\midrule
200 & 1 & 1 & 200 & 0.34 & baseline \\
200 & 3 & 1 & 600 & 0.42 & +0.08 \\
200 & 6 & 1 & 1,200 & 0.49 & +0.15 \\
200 & 6 & 6 & 7,200 & 0.54 & +0.20 \\
600 & 1 & 1 & 600 & 0.71 & +0.37 \\
600 & 3 & 1 & 1,800 & 0.82 & +0.48 \\
600 & 3 & 2 & 3,600 & \textbf{0.90} & +0.56 \\
1,000 & 3 & 2 & 6,000 & 0.93 & +0.59 \\
1,200 & 4 & 3 & 14,400 & 0.95 & +0.61 \\
\bottomrule
\end{tabular}
\end{table}

Several patterns emerge.
First, \textbf{current practice is severely under-powered}: $G = 0.34$ for the standard 1-prompt, 1-seed, 200-item configuration, far below the $G \geq 0.80$ threshold considered minimally acceptable in psychometrics.
Second, \textbf{item count dominates the budget}: tripling items from 200 to 600 (at 1 prompt, 1 seed) raises $G$ from 0.34 to 0.71---a gain of +0.37---while tripling prompts from 1 to 3 (at 200 items) raises $G$ only to 0.42 (+0.08).
The marginal effectiveness ratio is 10.2:1 (items vs.\ prompts) and 51.0:1 (items vs.\ seeds).
Third, $G \geq 0.90$ is achievable with 600 items, 3 prompts, and 2 seeds (3,600 total evaluations), representing an 18$\times$ cost factor over current practice.

Table~\ref{tab:cross_benchmark_budget} extends these projections across benchmarks.
Generation-heavy tasks require more items: GSM8K needs 800 items and MATH needs 1,100 items for $G \geq 0.90$, compared to 600 for MMLU and 550 for HellaSwag.

\begin{table}[t]
\centering
\caption{Minimum budget for $G \geq 0.90$ across benchmarks (3 prompts, 2 seeds).}
\label{tab:cross_benchmark_budget}
\small
\begin{tabular}{lrrrl}
\toprule
Benchmark & Items & Evals & Cost Factor \\
\midrule
HellaSwag & 550 & 3,300 & 16.5$\times$ \\
MMLU & 600 & 3,600 & 18$\times$ \\
ARC-Chall. & 650 & 3,900 & 19.5$\times$ \\
GSM8K & 800 & 4,800 & 24$\times$ \\
MATH & 1,100 & 13,200 & 66$\times$ \\
\bottomrule
\end{tabular}
\end{table}

\subsection{Scale Invariance: 14B Models}

Exp.~5 tests whether item dominance is an artifact of the 7--9B model scale by evaluating two 14B models.
Table~\ref{tab:14b_comparison} compares variance structures.

\begin{table}[t]
\centering
\caption{Variance pattern comparison: 8B models (mean and range) vs.\ 14B models on MMLU.}
\label{tab:14b_comparison}
\small
\begin{tabular}{lrrrr}
\toprule
Component & \makecell{8B \\ Mean} & \makecell{8B \\ Range} & \makecell{Qwen3 \\ 14B} & \makecell{DS-R1 \\ 14B} \\
\midrule
item\_id & 58.4 & 54.7--62.0 & 57.4 & 53.8 \\
residual & 16.8 & 14.9--18.5 & 16.2 & 19.8 \\
temp$\times$item & 5.8 & 4.6--7.0 & 5.3 & 8.2 \\
prompt$\times$item & 6.1 & 4.7--7.5 & 5.8 & 5.4 \\
precision & 0.3 & 0.1--0.5 & 0.2 & 0.2 \\
\bottomrule
\end{tabular}
\end{table}

Both 14B models show item\_id as the rank-1 variance component (57.4\% for Qwen3-14B, 53.8\% for DeepSeek-R1-Distill-14B), within the range observed for 8B models.
DeepSeek-R1-Distill-14B shows elevated residual (19.8\%) and temperature $\times$ item interaction (8.2\%), consistent with its reasoning-distilled architecture introducing additional chain-of-thought stochasticity.
Precision main effects remain negligible at 0.2\% for both models.
These results confirm that item dominance is not a small-model artifact.

\subsection{Fine-Tuning Variance Structure}

Exp.~6 extends the analysis from inference to fine-tuning by conducting a G-study of RoBERTa-large across 4 NLU tasks (SST-2, MNLI, RTE, COPA) with facets: training seed (6), learning rate (3), data shuffle order (4), GPU (3), and precision (2).
Table~\ref{tab:finetuning} compares the variance hierarchy between inference and fine-tuning.

\begin{table}[t]
\centering
\caption{Variance hierarchy comparison: inference (MMLU) vs.\ fine-tuning (SST-2).}
\label{tab:finetuning}
\small
\begin{tabular}{clcl}
\toprule
Rank & Inference (MMLU) & & Fine-tuning (SST-2) \\
\midrule
1 & item\_id (58.4\%) & & item\_id (41.2\%) \\
2 & residual (16.8\%) & & residual (13.1\%) \\
3 & prompt$\times$item (6.1\%) & & seed$\times$item (9.8\%) \\
4 & temp$\times$item (5.8\%) & & seed (8.7\%) \\
5 & seed$\times$item (2.9\%) & & lr$\times$item (7.4\%) \\
6 & prec$\times$item (1.8\%) & & lr (6.3\%) \\
\bottomrule
\end{tabular}
\end{table}

Item\_id remains the rank-1 component across all 4 fine-tuning tasks, ranging from 36.8\% (RTE) to 44.5\% (MNLI).
The key structural difference is that training seed and learning rate emerge as substantial variance sources (combined 10--22\%), replacing prompt and temperature as the second-tier factors.
RTE shows the lowest item\_id share (36.8\%) and highest seed main effect (12.4\%), consistent with its small training set (2.5K examples) making optimization more sensitive to initialization.
Infrastructure effects (GPU, precision) remain negligible ($<$0.5\% combined) in fine-tuning, and the current-practice $G$ for fine-tuning is 0.28---even lower than inference.

\subsection{Extended Precision Analysis}

Exp.~7 evaluates LayerCast (FP16 storage + FP32 compute) as an additional precision mode.
Table~\ref{tab:precision_modes} summarizes all precision modes.

\begin{table}[t]
\centering
\caption{Precision mode comparison on MMLU (4 models pooled).}
\label{tab:precision_modes}
\small
\begin{tabular}{lrrrr}
\toprule
Mode & Acc. & \makecell{Item Agr. \\ w/ FP32} & \makecell{Tput \\ (tok/s)} & \makecell{VRAM \\ (GB)} \\
\midrule
FP32 & .666 & 1.000 & 1,820 & 31.2 \\
BF16 & .671 & .999 & 2,760 & 15.8 \\
FP16 & .669 & .998 & 2,810 & 15.6 \\
LayerCast & .670 & .999 & 2,510 & 16.1 \\
\bottomrule
\end{tabular}
\end{table}

After the engine confound fix, all precision modes produce accuracy within a 0.5 percentage-point band.
LayerCast achieves item-level agreement of 0.9991 with FP32 while using 48\% less VRAM and delivering 38\% higher throughput, making it a practical option for practitioners seeking FP32-equivalent behavior at reduced cost.
LayerCast contributes only 0.04\% additional variance to the precision component.

\section{Discussion}

\paragraph{Why does item selection dominate?}
The magnitude of item\_id variance (42--58\%) reflects the steep difficulty gradient within every benchmark we studied.
MMLU items range from trivial factual recall to expert-level reasoning; GSM8K spans single-step arithmetic to multi-step word problems; MATH includes Level 1 algebra through Level 5 competition problems.
This difficulty gradient produces consistent accuracy ceilings and floors across all evaluation conditions, generating variance that dwarfs the effects of prompt, seed, or hardware changes.
The implication is fundamental: two evaluations using different 200-item subsets of the same benchmark may disagree more than two evaluations using different prompts, seeds, or GPUs on the same items.

\paragraph{Task difficulty modulates variance structure.}
The five benchmarks reveal a systematic gradient: as task difficulty increases (lower accuracy), item\_id share decreases (58.4\% $\to$ 42.1\%), residual increases (16.8\% $\to$ 24.8\%), and temperature $\times$ item interaction increases (5.8\% $\to$ 10.4\%).
This gradient quantifies why reasoning benchmark evaluations are less reproducible than classification benchmarks, and provides task-specific guidance for evaluation design.

\paragraph{Separability: a nuanced picture.}
The aggregate separability result ($<$0.5\% direct interactions) is good news for the many studies that isolate individual factors.
However, the substantial item-conditional interactions (14.2\% combined in Exp.~3) introduce an important caveat: conclusions from isolated-factor studies hold for the specific item sample used but may not generalize to different item samples.
This is not a flaw in existing studies but rather a boundary condition that should be acknowledged.

\paragraph{Implications for evaluation practice.}
Our D-study budgets translate directly into recommendations: for $G \geq 0.90$ reliability on MMLU, evaluate on 600 items with 3 prompt templates and 2 random seeds (3,600 evaluations total).
For harder benchmarks like MATH, the budget increases to 1,100 items.
The key insight is that item sampling provides $10\times$ the marginal return of prompt replication and $51\times$ the return of seed replication.
Practitioners constrained by compute should prioritize more items over more prompts or seeds.

\paragraph{Limitations.}
Our study is limited to open-weight models where infrastructure is controllable; API models (studied by \citealt{mesting2026}) have opaque infrastructure that may introduce additional variance.
We study 7--14B models; very large models ($>$70B) may exhibit different variance structures due to capacity effects.
The judge is fixed across experiments; adding judge as a facet (as in \citealt{mesting2026}) would likely increase total variance and could alter relative component shares.
Our benchmarks span multiple-choice, generation, and reasoning tasks but do not cover open-ended generation or safety evaluation.

\section{Conclusion}

We present the first fully-crossed G-theory study unifying infrastructure and methodology variance sources in LLM evaluation.
Three findings emerge consistently across 6 models, 5 benchmarks, and over 2 million evaluation records:
(1) Item selection is the dominant variance source (42--58\%), exceeding all other factors combined.
(2) Infrastructure and methodology factors are separable at the aggregate level (direct interaction $<$0.5\%) but couple through item difficulty (14.2\%), making conclusions from isolated-factor studies item-conditional.
(3) Current evaluation practices achieve $G < 0.40$; D-study budgets show that tripling item count is $10\times$ more cost-effective than tripling prompt count for reaching $G \geq 0.90$.
These findings reframe the reproducibility conversation: the field's focus on standardizing prompts, seeds, and hardware addresses secondary variance sources while leaving the largest source---item selection---uncontrolled.

\begin{thebibliography}{10}

\bibitem[Brittlebench(2026)]{brittlebench2026}
Brittlebench Authors.
\newblock Brittlebench: Quantifying prompt brittleness in LLM evaluation.
\newblock In \emph{Proceedings of NAACL}, 2026.

\bibitem[Cronbach et~al.(1972)]{cronbach1972}
L.~J. Cronbach, G.~C. Gleser, H.~Nanda, and N.~Rajaratnam.
\newblock \emph{The Dependability of Behavioral Measurements: Theory of Generalizability for Scores and Profiles}.
\newblock Wiley, 1972.

\bibitem[Hua et~al.(2025)]{hua2025}
W.~Hua et~al.
\newblock Are we evaluating LLMs correctly? Rethinking scoring with generation-based evaluation.
\newblock In \emph{Proceedings of EMNLP}, 2025.

\bibitem[Messina et~al.(2026)]{messina2026}
A.~Messina et~al.
\newblock Background temperature in commercial LLM APIs.
\newblock \emph{Transactions on Machine Learning Research}, 2026.

\bibitem[Mesting et~al.(2026)]{mesting2026}
B.~Mesting et~al.
\newblock Generalizability theory for LLM evaluation: Towards reliable model comparison.
\newblock In \emph{Proceedings of ACL}, 2026.

\bibitem[ThinkingMachines(2025)]{thinkingmachines2025}
ThinkingMachines.
\newblock 80 unique completions at temperature zero: Batch-invariance failures in LLM inference.
\newblock \emph{arXiv preprint arXiv:2501.xxxxx}, 2025.

\bibitem[Yuan et~al.(2025)]{yuan2025}
Z.~Yuan et~al.
\newblock Numerical precision matters: Variance in LLM inference across hardware configurations.
\newblock In \emph{Proceedings of NeurIPS}, 2025.

\end{thebibliography}

\end{document}
